% =====================================================================
%  CONJURA SOLUTION / PARTIAL PROOF
%  Extraction and decomposition for two split unpredictable sources.
%  Requires conjura-solution.cls and conjura-conjecture.cls here.
% =====================================================================
\documentclass{conjura-solution}

\runninghead{EXTRACTION FOR SPLIT UNPREDICTABLE SOURCES}

\newcommand{\Fun}{\mathsf{Fun}}
\newcommand{\Adv}{\mathsf{Adv}}
\newcommand{\Real}{\mathsf{Real}}
\newcommand{\Dec}{\mathsf{Dec}}
\newcommand{\Prob}[1]{\Pr\!\left[#1\right]}
\newcommand{\Exp}{\mathbb{E}}
\newcommand{\bits}{\{0,1\}}
\newcommand{\vx}{\mathbf{x}}
\newcommand{\vz}{\mathbf{z}}
\newcommand{\vu}{\mathbf{u}}
\newcommand{\ind}{\mathbb{1}}
\newcommand{\Unif}{\mathrm{Unif}}
\newcommand{\supp}{\mathrm{supp}}
\newcommand{\kna}{\kappa^{\mathrm{na}}}

\begin{document}

\cjkicker{CONJURA \textperiodcentered{} PARTIAL SOLUTION}
\cjtitle{Extraction for Split Unpredictable Sources, and Decomposition on a Region}
\cjresolves{\emph{Decomposition for Split Unpredictable Sources} (see
  \texttt{main.tex} in this folder). \textbf{Partially.} Two unconditional
  extraction bounds are proved outright; the conjecture itself is proved on a
  restricted region of $P$, $q$ and $M$, and remains open in general.}

\vspace{0.4em}

\section*{Reader's guide}

\noindent\textbf{What is proved.} Write $\kna(q)$ for the extraction advantage
restricted to observers whose query positions, for each fixed coin string, do not
depend on the challenge value, and $\mu(s):=\min\bigl(s\delta,\,3(s\delta^{2})^{1/3},\,1\bigr)$.

\begin{itemize}
\item[(A)] $\kna(q)\le 5\sqrt{\sigma'\delta}+q\delta$ for every $q$, and
  $\kappa(0)\le5\sqrt{\sigma'\delta}$ against \emph{all} observers.
  Unconditional (Theorem~\ref{thm:A}). This settles the statement's assertion that queries whose
  positions ignore the challenge value are almost free at every budget: their
  cost is the additive $q\delta$, with no multiplier on the leading term and no
  dependence on $M$.
\item[(B)] $\kappa(q)\le 5\sqrt{\sigma'\delta}+\mu(\min(qM,N^{2}))$ for every $q$
  and every observer. Unconditional (Theorem~\ref{thm:B}).
\item[(C)] An explicit family, chosen before $q$, indexed by the oracle and the
  leakage, and consistent with the oracle, with
  $\Adv_{\mathcal Y,D}\le\varepsilon(D)+2(\sigma'+\log\gamma^{-1})q/P+\gamma+P\delta+q\delta$ (Theorem~\ref{thm:C}).
\item[(C$'$)] Hence the conjecture, with the absolute constants $c=2$ and $C=8$, for
  every $P\le\sqrt{\sigma'/\delta}$ and every $q$ with $q^{+}\delta\le\sigma'$,
  provided either the observer's query positions are challenge-blind per coin
  string, or $M\le\sigma'/\sqrt{27\delta}$ (Theorem~\ref{thm:cprime}).
\end{itemize}

\noindent\textbf{What is not proved.} The conjecture in full. Two obstructions
survive, stated as mechanisms in \S\ref{sec:barriers}: challenge-steered probes
at large $M$, and large $P$ with growing $q$. The statement's own remark already
names the second as a separate problem.

\medskip

\noindent\textbf{Where to be most sceptical.} Two rounds of blind review found
two serious defects, both since repaired, and both of the same species: a
claim left \emph{underdetermined} rather than false. They were at
Theorem~\ref{thm:cprime}'s first proof line (which quantity the cited corollary
actually bounds) and in Definition~\ref{def:res} (whether challenge resolution is
a condition on the law of the query positions or on each coin string). Those two
places are where a third error, if there is one, is most likely to be. The review
status is recorded honestly in \S\ref{sec:review}.

\medskip

\noindent\textbf{Notation} is that of the statement document and is not repeated:
$\Fun=\{f:[N]\times[N]\to[M]\}$, $H\sample\Fun$, split $\delta$-unpredictable
sources $(S_1,S_2)$ with leakage bounded by $\sigma_1,\sigma_2$,
$\sigma=\sigma_1+\sigma_2$, $\sigma'=\sigma+2\log N$, $q^{+}=q+1$, and the
experiments $\Real,\Dec,\Real_0,\Dec_0$ with $\Adv_{\mathcal Y,D}$ and
$\kappa(q)$. Throughout $N\ge2$ and $M\ge2$; logarithms are base two. Write
$\mathcal Z:=\bits^{\le\sigma_1}\times\bits^{\le\sigma_2}$, so
$|\mathcal Z|=(2^{\sigma_1+1}-1)(2^{\sigma_2+1}-1)<2^{\sigma+2}$.

\section{The two structural facts}

Everything below rests on two observations, which together turn the conjecture
into a statement about the discrepancy of a random matrix on combinatorial
rectangles.

\begin{lemma}[derandomisation]\label{lem:derand}
For every $q$ and every $q$-query challenge-oblivious $D$ there is a
deterministic such $D'$ with \emph{(i)}
$\Adv_{\mathcal Y,D}\le\Adv_{\mathcal Y,D'}$ and
$\bigl|\Prob{\Real=1}-\Prob{\Real_0=1}\bigr|$ no larger for $D$ than for $D'$,
and \emph{(ii)} $D'$ of the same challenge resolution as $D$
(Definition~\ref{def:res}). Consequently the suprema defining $\kappa(q)$ and
$\kna(q)$ are both attained on deterministic observers, the latter on
deterministic observers of resolution~$1$.
\end{lemma}

\begin{proof}
Let $\rho$ be $D$'s coins and $D_\rho$ the deterministic strategy it then runs.
All of $\Fun,[N],[M],\mathcal Z$ are finite and $D$ makes at most $q$ queries, so
there are finitely many deterministic $q$-query challenge-oblivious strategies.
Each quantity at issue is $|\Exp_\rho[\alpha(\rho)]|$ with
$\alpha(\rho)=\Prob{\mathsf E_1=1\mid\rho}-\Prob{\mathsf E_0=1\mid\rho}$, the
coins being independent of every other draw and the paired experiments differing
in no other component; so it is at most $\max_\rho|\alpha(\rho)|$, attained at
some $\rho^{*}$. Take $D':=D_{\rho^{*}}$, which is $q$-query and
challenge-oblivious. It inherits $D$'s challenge resolution: the definition
quantifies over every coin string, so if $D^{f}_{\rho}(v,\zeta)$ and
$D^{f}_{\rho}(v',\zeta)$ issue the same query positions for every $\rho$ whenever
$h(v)=h(v')$, they do so in particular at $\rho^{*}$, which is what $D'$ runs.
The argument is identical for either paired experiment, giving both parts.
\end{proof}

Once $D$ is deterministic and its oracle is fixed to $f$, its output is a
function of its challenge input alone; write $\theta:[M]\to\bits$ for that
function and $p_\theta:=|\theta^{-1}(1)|/M$. That the test is \emph{named by
short data} is the crux of Lemma~\ref{lem:disc}.

\begin{lemma}[posterior product structure]\label{lem:prod}
For every $(f,\zeta)$ in the support of $(H,\vz)$ the conditional law of $\vx$ is
the product $\Pi_{f,\zeta}=\pi^{1}_{f,\zeta_1}\otimes\pi^{2}_{f,\zeta_2}$, where
$\pi^{i}_{f,\zeta_i}(u)=\Prob{x_i=u\mid H=f,\ z_i=\zeta_i}$. Writing
$m_i:=\|\pi^{i}_{f,\zeta_i}\|_\infty$:
\begin{enumerate}
\item $\max_{\vu\in[N]^2}\Pi_{f,\zeta}(\vu)=m_1m_2$;
\item $\Exp[m_i]\le\delta$ for $i=1,2$;
\item $\Exp[m_1m_2]\le\delta$;
\item $\Prob{m_1m_2>\theta}\le2\delta/\sqrt\theta$ for every $\theta\in(0,1]$.
\end{enumerate}
\end{lemma}

\begin{proof}
The sources run on independent private coins and on the same $H$, so conditioned
on $H=f$ their output pairs are independent:
$\Prob{x_1=u_1,z_1=\zeta_1,x_2=u_2,z_2=\zeta_2\mid f}=\prod_i\Prob{x_i=u_i,z_i=\zeta_i\mid f}$.
Summing over $u_1,u_2$ gives
$\Prob{\vz=\zeta\mid f}=\prod_i\Prob{z_i=\zeta_i\mid f}>0$, and dividing gives the
product form. Item~1 is the fact that the maximum of a product measure is the
product of the maxima.

For Item~2, let $\mathsf P$ be the predictor that, on input $\vz$ and with oracle
access to $H$, reads all of $H$, computes $\pi^{i}_{H,z_i}$ (it may, the
sources being fixed and predictors being arbitrary, possibly randomised,
functions of their inputs, as the statement document's own proof of its
fixed-set lemma instructs), and outputs a maximiser. Then
$\Prob{\mathsf P^{H}(\vz)=x_i}=\Exp_{(f,\zeta)}[m_i(f,\zeta)]$, which
unpredictability bounds by $\delta$. Item~3 follows from $m_2\le1$ and Item~2.
For Item~4, $m_1m_2>\theta$ with both factors in $[0,1]$ forces
$m_i>\sqrt\theta$ for some $i$; Markov and Item~2 give
$\Prob{m_i>\sqrt\theta}\le\delta/\sqrt\theta$.
\end{proof}

Item~1 is where the second source is spent: a \emph{single cell} of $[N]^2$
carries mass $m_1m_2$, a product of two independently small numbers, whereas one
unpredictable source over $[N]^2$ gives only $\delta$ per cell.

\begin{remark}[{$\Exp[m_1m_2]\le\delta^{2}$ is false}]\label{rem:corr}
The failure forces the exponent $1/3$ in
Lemma~\ref{lem:mass}. Let $2\le M\le N$, put $\delta:=1/M$, let $E$ be the event
$f(1,1)=1$, of probability $\delta$, and let both sources output $x_i:=1$ if $E$
holds and $x_i$ uniform on $[N]$ otherwise, leaking the empty string. The coins
are independent, so the pair is split, and
$\Exp[m_i]=\delta+(1-\delta)/N\le2\delta$, so the pair is
$2\delta$-unpredictable; yet $m_1m_2=1$ on $E$, so $\Exp[m_1m_2]\ge\delta$.
Splitness constrains the coins, not the two sources' common dependence on $f$.
\end{remark}

\section{Flattening and rectangle discrepancy}

\begin{lemma}[flattening]\label{lem:flat}
Let $\pi$ be a distribution on $[N]$ with $\|\pi\|_\infty=m$ and
$k:=\lfloor1/m\rfloor$. Then $1\le k\le N$, $1/k\le2m$, and $\pi$ is a convex
combination of distributions each uniform on a subset of size exactly $k$.
Consequently, for every $G:[N]^{2}\to[-1,1]$ and every $\pi^{1},\pi^{2}$ with
$\|\pi^{i}\|_\infty=m_i$, writing $k_i:=\lfloor1/m_i\rfloor$,
\[
  \bigl|\Exp_{\pi^{1}\otimes\pi^{2}}[G]\bigr|\ \le\
  \max\bigl\{\ \bigl|\Exp_{\Unif(B_1)\otimes\Unif(B_2)}[G]\bigr|\ :\ |B_i|=k_i\ \bigr\}.
\]
\end{lemma}

\begin{proof}
The inequality $m\ge1/N$ gives $k\le N$; $m\le1$ gives $k\ge1$. For $x\ge1$ we
have $\lfloor x\rfloor\ge x/2$ (for $1\le x<2$ because $\lfloor x\rfloor=1\ge x/2$;
for $x\ge2$ because $\lfloor x\rfloor>x-1\ge x/2$), so taking $x=1/m$ yields
$1/k\le2m$. If $k=1$ the claimed decomposition is the identity
$\pi=\sum_u\pi(u)\delta_u$, a convex combination of distributions uniform on
singletons, and no external result is needed; assume $k\ge2$. Since
$\|\pi\|_\infty\le1/k=2^{-\log k}$ with $\log k>0$ and $2^{\log k}=k\in\mathbb N$,
the flat-decomposition lemma [CFHS, Lemma~3] applies and writes $\pi$ as a convex
combination of distributions uniform on $k$-subsets. The displayed inequality
follows because $(\rho^1,\rho^2)\mapsto\Exp_{\rho^1\otimes\rho^2}[G]$ is
bilinear, so substituting both decompositions expresses the left side as a convex
combination of the quantities on the right.
\end{proof}

Note $G$ is arbitrary: it will carry an indicator $\ind[\vu\notin\mathcal S]$ that
does \emph{not} factor over the coordinates. The failure to factor is harmless, because the linear
structure used lives in $(\pi^1,\pi^2)$, not in $G$.

\begin{definition}[revealing rule]\label{def:rr}
A \emph{revealing rule of budget $s$} is a family
$(\mathcal A_\zeta)_{\zeta\in\mathcal Z}$ of procedures, where $\mathcal A_\zeta$
inspects the coordinates of $f$ one at a time (the next a function of $\zeta$
and the values already seen), halts having inspected a set
$\mathcal S(f,\zeta)$ with $|\mathcal S(f,\zeta)|\le s$, and outputs a test
$\theta_{\zeta,f}:[M]\to\bits$ that is a function of $\zeta$ and the inspected
values alone.
\end{definition}

\begin{lemma}[rectangle discrepancy against leakage-named tests]\label{lem:disc}
Let $\gamma_0\in(0,1)$ and let $(\mathcal A_\zeta)$ be a revealing rule. Put
$L:=\ln(eN)$, $C_0:=(\sigma+2)\ln2+\ln\frac{4N^{2}}{\gamma_0}$, and
\[
  t(k_1,k_2):=\sqrt{\frac{k_1\ln\frac{eN}{k_1}+k_2\ln\frac{eN}{k_2}+C_0}{2k_1k_2}}\,.
\]
Then with probability at least $1-\gamma_0$ over $f\sample\Fun$, simultaneously
for every $\zeta\in\mathcal Z$, all $k_1,k_2\in[N]$ and all $B_i\subseteq[N]$
with $|B_i|=k_i$, writing $R:=B_1\times B_2$,
\begin{equation}\label{eq:disc}
  \Bigl|\ \frac{1}{k_1k_2}\!\!\sum_{u\in R\setminus\mathcal S(f,\zeta)}\!\!
  \bigl(\theta_{\zeta,f}(f(u))-p_{\theta_{\zeta,f}}\bigr)\Bigr|\ \le\ t(k_1,k_2).
\end{equation}
Moreover $t(k_1,k_2)\le\sqrt{L(m_1+m_2)+2C_0m_1m_2}$ whenever
$k_i=\lfloor1/m_i\rfloor$, and hence
\begin{equation}\label{eq:Et}
  \Exp\bigl[t(k_1,k_2)\bigr]\ \le\ \sqrt{2\delta(L+C_0)}\ \le\
  \sqrt{8\delta\bigl(\sigma'+\log\gamma_0^{-1}\bigr)}.
\end{equation}
\end{lemma}

\begin{proof}
\emph{Conditional uniformity off the revealed set.} Fix $\zeta$. For
$S_0\subseteq[N]^{2}$ and $w\in[M]^{S_0}$ the event
$\mathcal T_{S_0,w}:=\{\mathcal S(f,\zeta)=S_0\wedge f|_{S_0}=w\}$ is measurable
with respect to $f|_{S_0}$: replaying $\mathcal A_\zeta$ on $w$ reproduces its
whole inspection sequence. The coordinates of $f$ being i.i.d.\ uniform,
conditioning on such an event leaves $f|_{[N]^{2}\setminus S_0}$ i.i.d.\ uniform;
and $\theta_{\zeta,f}$, $p_{\theta_{\zeta,f}}$ are constant on
$\mathcal T_{S_0,w}$.

\emph{One tuple.} Fix also $k_1,k_2,B_1,B_2$ and condition on
$\mathcal T_{S_0,w}$. Write $\theta,p_\theta$ for the constants and
$n_0:=|R\setminus S_0|\le k_1k_2$. The variables $\theta(f(u))-p_\theta$,
for $u\in R\setminus S_0$, are independent, lie in an interval of length $1$, and have
mean $0$, since $f(u)$ is uniform on $[M]$ and
$\Prob{\theta(f(u))=1}=p_\theta$ exactly. Hoeffding's inequality with
$\lambda:=k_1k_2\,t(k_1,k_2)$ gives conditional failure probability at most
$2\exp(-2\lambda^{2}/n_0)\le2\exp(-2t^{2}k_1k_2)$, uniformly in $(S_0,w)$, hence
also unconditionally.

\emph{Union bound.} By the definition of $t$,
$2\exp(-2t^{2}k_1k_2)=2\,e^{-k_1\ln\frac{eN}{k_1}}e^{-k_2\ln\frac{eN}{k_2}}2^{-(\sigma+2)}\frac{\gamma_0}{4N^{2}}$.
There are $|\mathcal Z|\binom{N}{k_1}\binom{N}{k_2}$ tuples with the given
$(k_1,k_2)$, and $\binom{N}{k}\le(eN/k)^{k}$ [CFHS, Lemma~4.3] while
$|\mathcal Z|<2^{\sigma+2}$. Multiplying, all three exponential factors cancel
and the contribution of the pair $(k_1,k_2)$ is at most
$2\gamma_0/(4N^{2})=\gamma_0/(2N^{2})$; summing over the $N^{2}$ pairs gives
$\gamma_0/2\le\gamma_0$.

\emph{Bounding $t$.} Expanding and using $k_i\ge1$, so
$\ln(eN/k_i)\le L$,
\[
  t(k_1,k_2)^2=\frac{\ln\frac{eN}{k_1}}{2k_2}+\frac{\ln\frac{eN}{k_2}}{2k_1}
  +\frac{C_0}{2k_1k_2}\ \le\ \frac{L}{2k_2}+\frac{L}{2k_1}+\frac{C_0}{2k_1k_2},
\]
and $1/k_i\le2m_i$ from Lemma~\ref{lem:flat} bounds the three terms by $Lm_2$,
$Lm_1$ and $2C_0m_1m_2$.

\emph{Expectations.} By concavity of $\sqrt{\cdot}$ and
Lemma~\ref{lem:prod}(2),(3),
$\Exp[t]\le\sqrt{L(\Exp m_1+\Exp m_2)+2C_0\Exp[m_1m_2]}\le\sqrt{2\delta(L+C_0)}$.
Numerically
$L+C_0=1+\ln N+(\sigma+2)\ln2+\ln4+2\ln N+\ln\gamma_0^{-1}\le3\ln N+0.694\sigma+3.78+\ln\gamma_0^{-1}$;
as $\sigma'\ge2.885\ln N$ and $\sigma'\ge2$ for $N\ge2$, this is at most
$3.63\sigma'+\ln\gamma_0^{-1}$, giving~\eqref{eq:Et}.
\end{proof}

\begin{remark}[why $M$ is absent, and where one source would break this]\label{rem:one}
The alphabet enters $t$ only through $p_\theta$, which cancels, and through the
i.i.d.\ structure of $\{f(u)\}$, which is alphabet-blind. There are $2^{M}$
possible tests, but the union runs over \emph{which} test, and once the observer
is fixed and the revealed values are conditioned away only $|\mathcal Z|<2^{\sigma+2}$
are reached. The narrow union is the one respect in which the route beats a Fourier or
chi-squared analysis, which would pay $\sqrt M$ to reduce an arbitrary test to
characters.

The statement document demands that every proof fail visibly when the two sources
are replaced by one. It does, and not in the size of $t$: at $k_1=1,k_2=N$ one
gets $t(1,N)^{2}=(L+N+C_0)/(2N)\to\tfrac12$, so $t\to1/\sqrt2$, neither large nor
vacuous; and in any case $k_1=1$ models a \emph{deterministic first
coordinate}, not a single source. The failure is earlier and total. A single
source emits a point of $[N]^{2}$ whose posterior given the oracle and the
leakage need not factor at all; Lemma~\ref{lem:prod}(1) is then unavailable, a
cell carries mass up to $\delta$ rather than $m_1m_2$, and there is no product
measure for Lemma~\ref{lem:flat} to flatten onto a rectangle. Lemmas~\ref{lem:flat}
and~\ref{lem:disc} have no input and the argument does not begin.
\end{remark}

\section{The observer as a revealing rule}

\begin{definition}[challenge resolution]\label{def:res}
An observer $D$ has \emph{challenge resolution $M'$} if there is $h:[M]\to[M']$ such that for
\textbf{every coin string $\rho$ of $D$}, every $f$, every $\zeta$, and all
$v,v'$ with $h(v)=h(v')$, the runs $D^{f}_{\rho}(v,\zeta)$ and
$D^{f}_{\rho}(v',\zeta)$ issue the same sequence of query \emph{positions}; the
output bit is unconstrained. Every $D$ has resolution $M$; resolution $1$ means
that, for each fixed coin string, the query positions ignore the challenge value.
We may take $h$ surjective without loss: replacing $[M']$ by $\mathrm{image}(h)$
issues fewer runs below and only shrinks the budget.
\end{definition}

The coin quantifier is part of the definition, not a technicality. The weaker
condition that the query positions be challenge-independent \emph{in law} defines
a strictly larger class: at $N=M=2$ and $q=1$, the observer that flips $\rho$ and
queries $(1,1)$ when $v\oplus\rho=1$ and $(2,2)$ otherwise has query position
uniform on $\{(1,1),(2,2)\}$ for both challenge values, yet each of its two coin
fixings has resolution $2$. Such an observer steers its probes by the challenge
once its coins are fixed (exactly the behaviour Barrier~1 leaves open), and
is not in the class $\kna$ ranges over.

\begin{lemma}[the observer as a revealing rule]\label{lem:rr}
Let $D$ be deterministic, $q$-query, challenge-oblivious, of resolution $M'$.
Then the procedures ``$\mathcal A_\zeta$: for $j=1,\dots,M'$ in order, run
$D^{f}(v_j,\zeta)$ for a fixed $v_j\in h^{-1}(j)$ (which exists, $h$ being
surjective), inspecting each queried cell not already known'' form a revealing
rule of budget $s=\min(qM',N^{2})$, with $\theta_{\zeta,f}(v):=D^{f}(v,\zeta)$.
\end{lemma}

\begin{proof}
The procedure $\mathcal A_\zeta$ inspects one coordinate at a time, the next being a function
of $\zeta$ and the values already seen: the loop and the representatives are
fixed in advance, and inside run $j$ the next position is a function of
$(v_j,\zeta)$ and previous answers. Each of the $M'$ runs issues at most $q$
positions, so at most $qM'$ cells, and at most $N^{2}$ exist. Finally, for
arbitrary $v$ with $j=h(v)$, the run $D^{f}(v,\zeta)$ issues exactly the
positions run $j$ issued, all inspected and recorded, so its whole execution,
hence its output bit, is reconstructible from
$(v,\zeta,\mathcal S,f|_{\mathcal S})$. Therefore the entire function
$\theta_{\zeta,f}$, and $p_{\theta_{\zeta,f}}$, is determined by
$(\zeta,\mathcal S,f|_{\mathcal S})$.
\end{proof}

The resolution parameter is the whole difference between the settled and the open
case. Lemma~\ref{lem:disc} needs the test \emph{as a function on all of} $[M]$,
because the reference $p_\theta$ averages over the alphabet; so every value the
observer might receive must have its query path revealed in advance, and an
observer that steers by $v$ forces $M'=M$.

\begin{lemma}[product mass of the revealed set]\label{lem:mass}
For a revealing rule of budget $s\ge1$,
$\Exp\bigl[\min(\Pi_{f,\zeta}(\mathcal S(f,\zeta)),1)\bigr]\le
\mu(s):=\min\bigl(s\delta,\ 3(s\delta^{2})^{1/3},\ 1\bigr)$; for $s=0$ the left
side is $0$.
\end{lemma}

\begin{proof}
By Lemma~\ref{lem:prod}(1),
$\Pi(\mathcal S)=\sum_{\vu\in\mathcal S}\pi^1(u_1)\pi^2(u_2)\le s\,m_1m_2$.
Taking expectations and using Lemma~\ref{lem:prod}(3) gives the arm $s\delta$;
the arm $1$ is trivial. For the third, fix $\theta\in(0,1]$, bound by $s\theta$ on
$\{m_1m_2\le\theta\}$ and by $1$ off it, and use Lemma~\ref{lem:prod}(4):
$\Exp[\min(\Pi(\mathcal S),1)]\le g(\theta):=s\theta+2\delta\theta^{-1/2}$. Then
$g'(\theta)=s-\delta\theta^{-3/2}$ vanishes at $\theta^{*}=(\delta/s)^{2/3}$,
which is $\le1$ since $s\ge1\ge\delta$, and $g''>0$; and
$g(\theta^{*})=s^{1/3}\delta^{2/3}+2s^{1/3}\delta^{2/3}=3(s\delta^{2})^{1/3}$.
\end{proof}

By Remark~\ref{rem:corr} the bound $s\delta^{2}$ is \emph{not} available, and the
exponent $1/3$ is exactly what the $\sqrt\theta$ of Lemma~\ref{lem:prod}(4)
produces. Lemma~\ref{lem:mass} is tight at $s=1$ for that example.

\section{Extraction}

\begin{lemma}\label{lem:ident}
For deterministic $D$ with associated $\theta_{\zeta,f}$,
$\Prob{\Real=1}=\Exp_{(f,\zeta)}\bigl[\Exp_{\vx\sim\Pi_{f,\zeta}}[\theta_{\zeta,f}(f(\vx))]\bigr]$
and $\Prob{\Real_0=1}=\Exp_{(f,\zeta)}[p_{\theta_{\zeta,f}}]$.
\end{lemma}

\begin{proof}
Condition on $(H,\vz)=(f,\zeta)$. In both experiments $D$'s oracle is $f$, so its
output on challenge $v$ is $\theta_{\zeta,f}(v)$. In $\Real$ the challenge is
$f(\vx)$ with $\vx\sim\Pi_{f,\zeta}$ by Lemma~\ref{lem:prod}; in $\Real_0$ it is
uniform and independent, with $\Exp_{v\sample[M]}[\theta(v)]=p_\theta$.
\end{proof}

So $\Real$ and $\Real_0$ differ in the challenge alone, the oracle being the true
$H$ in both: the whole problem is the one-cell question \emph{is $f(\vx)$
distributed like a fresh uniform value, to a reader who knows $\zeta$ and can
probe $f$ in $q$ queries?}

\begin{proposition}\label{prop:master}
For every $\gamma_0\in(0,1)$ and every deterministic $q$-query
challenge-oblivious $D$ of resolution $M'$,
\[
  \bigl|\Prob{\Real=1}-\Prob{\Real_0=1}\bigr|\le\gamma_0
  +\sqrt{8\delta(\sigma'+\log\gamma_0^{-1})}+\mu\bigl(\min(qM',N^{2})\bigr).
\]
\end{proposition}

\begin{proof}
By Lemma~\ref{lem:ident} and the triangle inequality the left side is at most
$\Exp[\Delta]$ with
$\Delta(f,\zeta):=|\Exp_{\Pi}[\theta(f(\vx))-p_\theta]|$. Put
$G_{f,\zeta}(\vu):=(\theta(f(\vu))-p_\theta)\ind[\vu\notin\mathcal S(f,\zeta)]\in[-1,1]$,
so that, since $|\theta-p_\theta|\le1$,
\[
  \Delta(f,\zeta)\ \le\ \bigl|\Exp_{\Pi_{f,\zeta}}[G_{f,\zeta}]\bigr|
  +\Pi_{f,\zeta}(\mathcal S(f,\zeta)),
\]
the two terms being an exhaustive split of the challenge cell. Lemma~\ref{lem:flat}
bounds the first by the largest $|\Exp_{\Unif(B_1)\otimes\Unif(B_2)}[G]|$ over
$|B_i|=k_i:=\lfloor1/m_i\rfloor$, and for such a rectangle that quantity is
\[
  \frac{1}{k_1k_2}\sum_{\vu\in R}G_{f,\zeta}(\vu)
  =\frac{1}{k_1k_2}\sum_{\vu\in R\setminus\mathcal S(f,\zeta)}\bigl(\theta(f(\vu))-p_\theta\bigr),
\]
verbatim the quantity Lemma~\ref{lem:disc} controls; on the event $\mathcal E$ of
that lemma it is at most $t(k_1,k_2)$. As $\Delta\le1$ always, splitting on
$\mathcal E$ gives
$\Exp[\Delta]\le\Prob{\mathcal E^{c}}+\Exp[t]+\Exp[\min(\Pi(\mathcal S),1)]$, all
terms nonnegative so the indicators may be dropped; these are at most
$\gamma_0$, \eqref{eq:Et}, and $\mu(s)$ by Lemmas~\ref{lem:rr} and~\ref{lem:mass}.
\end{proof}

\begin{theorem}[extraction, challenge-blind query positions]\label{thm:A}
$\kna(q)\le5\sqrt{\sigma'\delta}+q\delta$, and
$\kappa(0)\le5\sqrt{\sigma'\delta}$ with no restriction on the observer.
\end{theorem}

\begin{proof}
Both $\kna(q)$ and $\kappa(0)$ are suprema over observers the statement permits
to randomise, whereas Proposition~\ref{prop:master} is stated for deterministic
$D$; by Lemma~\ref{lem:derand} the suprema are attained on deterministic
observers and $D'$ inherits $D$'s challenge resolution, so restricting to
deterministic resolution-$1$ observers loses nothing. If $\delta=1$ then
$\sqrt{\sigma'\delta}\ge\sqrt2>1$ and there is nothing to prove; else take
$\gamma_0:=\delta$ in Proposition~\ref{prop:master}. Since $\delta\ge1/N$ we get
$\log\gamma_0^{-1}\le\log N\le\sigma'/2$, so the middle term is at most
$\sqrt{12\sigma'\delta}\le3.47\sqrt{\sigma'\delta}$, and
$\gamma_0=\delta\le\sqrt{\sigma'\delta}$ because $\delta\le1\le\sigma'$.
Resolution $1$ gives $s\le q$ and $\mu(q)\le q\delta$. For $q=0$ we have $s=0$, and every
observer has resolution $1$ vacuously.
\end{proof}

\begin{corollary}\label{cor:A}
If $q^{+}\delta\le\sigma'$ then $\kna(q)\le6\sqrt{\sigma'q^{+}\delta}$. Outside
that regime $q^{+}\delta>\sigma'\ge2$, so $\sqrt{\sigma'q^{+}\delta}>\sigma'>1$
and the target bound is vacuous.
\end{corollary}

\begin{proof}
Monotonicity gives $5\sqrt{\sigma'\delta}\le5\sqrt{\sigma'q^{+}\delta}$; and $q^{+}\delta\le\sigma'$
gives $(q^{+}\delta)^{2}\le\sigma'q^{+}\delta$, i.e.
$q\delta\le q^{+}\delta\le\sqrt{\sigma'q^{+}\delta}$.
\end{proof}

The corollary settles the statement's line that queries whose positions ignore the
challenge value are almost free at every budget: their cost is the additive
$q\delta$, with no multiplier on the leading term and no $\log M$.

\begin{theorem}[extraction, unrestricted observers]\label{thm:B}
$\kappa(q)\le5\sqrt{\sigma'\delta}+\mu(\min(qM,N^{2}))$ for every $q$.
\end{theorem}

\begin{proof}
By Lemma~\ref{lem:derand} the supremum is attained on deterministic observers;
apply Proposition~\ref{prop:master} to such a $D$ with $M'=M$ and
$\gamma_0=\delta$, as in Theorem~\ref{thm:A}.
\end{proof}

\begin{corollary}\label{cor:B}
$\kappa(q)\le6\sqrt{\sigma'q^{+}\delta}$ whenever $q^{+}\delta\le\sigma'$ and
either $M\le\sqrt{\sigma'/(q\delta)}$ or
$M\le\sigma'^{3/2}\sqrt{q/\delta}/27$. If $M\le\sigma'/\sqrt{27\delta}$ then one
of the two holds for every $q\ge1$.
\end{corollary}

\begin{proof}
The inequality $qM\delta\le\sqrt{\sigma'q^{+}\delta}$ follows from $qM^{2}\delta\le\sigma'$, and
$3(qM\delta^{2})^{1/3}\le\sqrt{\sigma'q^{+}\delta}$ from
$27qM\delta^{1/2}\le\sigma'^{3/2}q^{3/2}$. The first arm fails only if
$q>\sigma'/(M^{2}\delta)$ and the second only if $q<729M^{2}\delta/\sigma'^{3}$;
both fail at some $q$ only if $\sigma'^{4}<729M^{4}\delta^{2}$, i.e.
$M>\sigma'/\sqrt{27\delta}$.
\end{proof}

\section{Presampling, and the decomposition}

\begin{lemma}[presampling, oracle-indexed and consistent]\label{lem:P}
Let $P\in\mathbb N$ and $\gamma\in(0,1)$. \emph{(For $M=1$ the set $\Fun$ is a singleton,
so $\Real_0$ and $\Dec_0$ are the same experiment and the bound holds with left
side $0$; the construction divides by $\log M$ and is stated for $M\ge2$.)} There
is a family $\mathcal Y^{P,\gamma}=\{Y_{f,\zeta}\}$ indexed by $f\in\Fun$ and
$\zeta\in\bits^{*}\times\bits^{*}$ (with $Y_{f,\zeta}$ uniform on $\Fun$ at
every $\zeta$ outside the support of $\vz$, so the family is total, no experiment
reaching such a $\zeta$), depending only on $(S_1,S_2,P,\gamma)$, in which every
$Y_{f,\zeta}$ is a $P$-mixture consistent with $f$, indeed a single
$P$-bit-fixing source, such that for every $q$ and every $q$-query observer
taking $\vz$ and an independent uniform value of $[M]$,
\[
  \bigl|\Prob{\Real_0=1}-\Prob{\Dec_0=1}\bigr|\ \le\
  \frac{q(\sigma+2+2\log\gamma^{-1})}{P}+2\gamma .
\]
\end{lemma}

\begin{proof}
\emph{Step 1: deficiency under randomised leakage.} For $\zeta$ with
$\Prob{\vz=\zeta}>0$ let $\mu_\zeta$ be the law of $H$ given $\vz=\zeta$ and
$S_\zeta:=N^{2}\log M-H_\infty(\mu_\zeta)$. Since
$\mu_\zeta(f)=\Prob{H=f}\Prob{\vz=\zeta\mid H=f}/\Prob{\vz=\zeta}\le(|\Fun|\Prob{\vz=\zeta})^{-1}$,
we get $S_\zeta\le\log(1/\Prob{\vz=\zeta})$. With
$\bar S:=\sigma+2+\log\gamma^{-1}$ and $\mathcal B:=\{\zeta:S_\zeta>\bar S\}$,
every $\zeta\in\mathcal B$ has $\Prob{\vz=\zeta}<\gamma2^{-(\sigma+2)}$, so
$\Prob{\vz\in\mathcal B}<|\mathcal Z|\gamma2^{-(\sigma+2)}<\gamma$. This replaces
the source's own expectation bound on the deficiency, which assumes the advice is
a deterministic function of the oracle and fails here.

\emph{Step 2: the family.} Fix $\zeta$ and set
$\delta_\zeta:=(S_\zeta+\log\gamma^{-1})/(P\log M)$. If $\delta_\zeta>1$ then
$(S_\zeta+\log\gamma^{-1})/P>\log M\ge1$, so the asserted bound exceeds
$q\ge1$ and is vacuous for $q\ge1$, while for $q=0$ both experiments give the
observer a uniform challenge and no oracle access; set $Y_{f,\zeta}$ uniform.
(The governing condition $\delta_\zeta>1$ is $q$-free, so the family is still
defined before $q$.) Otherwise apply the decomposition [CDGS, Claim~2] to
$\mu_\zeta$ with parameter $\delta_\zeta$; since
$(S_\zeta+\log\gamma^{-1})/(\delta_\zeta\log M)=P$ it yields
$\mu_\zeta=\sum_j\lambda_jX_j+\lambda_{\mathrm{fin}}Y_{\mathrm{fin}}$ with
$\lambda_{\mathrm{fin}}\le\gamma$, each $X_j$ a $(P,1-\delta_\zeta)$-dense
source, all supports pairwise disjoint, and each $X_j$ fixing its set $I_j$,
with $|I_j|\le P$, to values that every $f$ in its support takes. Let $Y_j$ be the
corresponding $P$-bit-fixing source and define $Y_{f,\zeta}:=Y_j$ for the unique
$j$ with $f\in\supp X_j$, and uniform otherwise. Consistency is immediate; the
construction reads only $(f,\zeta)$ and the decomposition, which depends only on
$(S_1,S_2,P,\gamma)$. Write $J=J(f,\zeta)$.

\emph{Step 3: one component.} Fix $\zeta\notin\mathcal B$ with $\delta_\zeta\le1$
and condition on $J=j$. Because the decomposition is of $\mu_\zeta$ itself,
$\lambda_jX_j$ is the restriction of $\mu_\zeta$ to $\supp X_j$, so the
conditional law of $H$ is exactly $X_j$; in $\Dec_0$ the oracle is drawn freshly
from $Y_j$. The challenge is uniform and independent, so conditioning on it
leaves a fixed $q$-query adaptive distinguisher between $X_j$ and $Y_j$;
[CDGS, Claim~3] bounds its advantage by
$q\delta_\zeta\log M=q(S_\zeta+\log\gamma^{-1})/P$, and averaging over the
challenge preserves this.

\emph{Step 4: averaging.} Averaging over $j$ and bounding the residue event, of
conditional probability at most $\gamma$, by $1$; then averaging over $\zeta$,
bounding $\mathcal B$'s contribution by $\Prob{\vz\in\mathcal B}<\gamma$ and
using $S_\zeta\le\bar S$ off $\mathcal B$, gives the claim.
\end{proof}

\begin{theorem}\label{thm:C}
With $\mathcal Y:=\mathcal Y^{P,\gamma/2}$, for every $q$ and every $q$-query
challenge-oblivious $D$, writing
$\varepsilon(D):=\bigl|\Prob{\Real=1}-\Prob{\Real_0=1}\bigr|$ for that same $D$,
\[
  \Adv_{\mathcal Y,D}\ \le\ \varepsilon(D)
  +\frac{2(\sigma'+\log\gamma^{-1})q}{P}+\gamma+P\delta+q\delta .
\]
Since $\varepsilon(D)\le\kappa(q)$, the same holds with $\kappa(q)$ in place of
$\varepsilon(D)$; but the $D$-relative form is what Theorem~\ref{thm:cprime}
uses, and the inequality $\kappa(q)\le\kna(q)$ is false in general and nowhere
needed.
\end{theorem}

\begin{proof}
Interpolate $\mathsf G_0=\Real$, $\mathsf G_1=\Real_0$, $\mathsf G_2=\Dec_0$ and
$\mathsf G_3=\Dec$.

$|\mathsf G_0-\mathsf G_1|=\varepsilon(D)$ by definition, for the fixed $D$ the
theorem is instantiated at; no supremum is taken here.

$|\mathsf G_1-\mathsf G_2|$: Lemma~\ref{lem:P} with slack $\gamma/2$ gives
$q(\sigma+2+2\log(2/\gamma))/P+\gamma$; as $N\ge2$ gives $\sigma+2\le\sigma'$ and
$\sigma'\ge2$, the numerator is at most $2\sigma'+2\log\gamma^{-1}$. For $q\ge P$
that term already exceeds $2\sigma'>1$, so no hypothesis $q<P$ is needed.

$|\mathsf G_2-\mathsf G_3|\le P\delta+q\delta$: draw $H$, $(\vx,\vz)$ and $J$
with fixed set $I_J$; let $H^{\circ}\sim Y_J$ and $U\sample[M]$ be independent of
everything else; let $H^{*}$ be $H^{\circ}$ with its value at $\vx$ overwritten by
$U$ when $\vx\notin I_J$. As $J$ is a function of $(H,\vz)$, the oracle $H^{\circ}$ is
independent of $\vx$ given $(H,\vz,J)$; conditioning on $\vx=\vu\notin I_J$, the
source $Y_J$ is uniform and independent at $\vu$, so both $H^{\circ}$ and
$H^{*}$ are distributed as $Y_J$. Running $\mathsf G_3$ with oracle $H^{*}$ and
$\mathsf G_2$ with oracle $H^{\circ}$ and challenge $U$, the two executions
receive the same input on $\{\vx\notin I_J\}$, since there $H^{*}(\vx)=U$, and
their oracles agree off $\vx$. So they diverge only if $\vx\in I_J$ or $D$ queries
$\vx$; being identical up to the step before a first query at $\vx$, that event
has equal probability in both and may be measured in $\mathsf G_2$. The statement
document's lemma \emph{the fixed set misses the challenge} gives
$\Prob{\vx\in I_J}\le P\delta$; its proof requires only that $\mathcal Y$
depend on $(S_1,S_2,P,\gamma)$ alone, be indexed by $(f,\zeta)$, and have its
component drawn independently of $\vx$ given $(H,\vz)$, all of which hold, $J$
being a function of $(H,\vz)$. Its lemma \emph{the observer misses the challenge}
gives $\Prob{\vx\in Q}\le q\delta$ in $\mathsf G_2$, its predictor being
simulable there because the challenge is uniform and independent.
\end{proof}

\begin{theorem}[the conjecture, on the region reached]\label{thm:cprime}
Suppose $q^{+}\delta\le\sigma'$, $P\le\sqrt{\sigma'/\delta}$, and either $D$ has
challenge resolution $1$, or the hypothesis of Corollary~\ref{cor:B} holds (in
particular if $M\le\sigma'/\sqrt{27\delta}$). Then with
$\mathcal Y=\mathcal Y^{P,\gamma/2}$,
\[
  \Adv_{\mathcal Y,D}\ \le\
  \frac{2(\sigma'+\log\gamma^{-1})q^{+}}{P}+8\sqrt{\sigma'q^{+}\delta}+\gamma,
\]
that is, the conjecture holds on that region with $c=2$ and $C=8$.
\end{theorem}

\begin{proof}
We have $\varepsilon(D)\le6\sqrt{\sigma'q^{+}\delta}$: in the first branch $D$ has
challenge resolution $1$, so $\varepsilon(D)\le\kna(q)$ and
Corollary~\ref{cor:A} applies (equivalently, Proposition~\ref{prop:master}
bounds $\varepsilon(D)$ directly at $M'=1$); in the second branch
Corollary~\ref{cor:B} bounds $\kappa(q)\ge\varepsilon(D)$. Then
$P\delta\le\sqrt{\sigma'\delta}\le\sqrt{\sigma'q^{+}\delta}$ by the hypothesis on
$P$, and $q\delta\le q^{+}\delta\le\sqrt{\sigma'q^{+}\delta}$ by
$q^{+}\delta\le\sigma'$. Summing, $6+1+1=8$.
\end{proof}

The family is built from $(P,\gamma)$ before $q$ and the bound holds for every
$q$, so the statement's quantifier order is respected; and every $Y_{f,\zeta}$ is
consistent with $f$, so the conclusion is the conjecture and not the weakening
its remark on consistency describes.

\section{Combination with the compression bound}

For $q\ge1$, [CFHS, Theorem~3] with $\ell=2$, $2^{-k}=\delta$ and $q_D=q$ gives
$\kappa(q)=O\bigl(\log M\cdot(q\delta^{2}(\sigma+2N))^{1/3}\bigr)$. Hence
\[
  \kappa(q)\le5\sqrt{\sigma'\delta}+\min\Bigl\{\mu(\min(qM,N^{2})),\
  O\bigl(\log M\,(q\delta^{2}(\sigma+2N))^{1/3}\bigr)\Bigr\},
\]
and Theorem~\ref{thm:C} converts either arm. The arms are complementary:
Theorem~\ref{thm:B} is smaller when $M\lesssim N\log^{3}M$, the compression arm
when $M$ is very large relative to $N$; the latter also relaxes as $q$ grows,
giving the target shape with no condition on $M$ once
$q\gtrsim\log^{6}M(\sigma+2N)^{2}\delta/\sigma'^{3}$. What neither arm reaches is
$M$ large together with $q$ small. Recorded against that source: its right-hand
side vanishes at $q_D=0$ whereas $\kappa(0)>0$ in general, so it is invoked only
for $q\ge1$, and its constant is not extracted, so this arm is asymptotic while
everything else here is explicit.

\section{What is not proved}\label{sec:barriers}

\paragraph{Barrier 1: challenge-steered probes.} Uncovered:
$M>\sigma'/\sqrt{27\delta}$ with $q$ in the window left by
Corollary~\ref{cor:B} and the compression arm. The gap is not slack in the counting.
Lemma~\ref{lem:disc} compares $\Exp_\Pi[\theta(f(\vx))]$ with $p_\theta$, an
average over \emph{all} of $[M]$, so $\theta$ must be known everywhere; by
Lemma~\ref{lem:rr} an observer that steers its probes by the challenge makes
$\theta(v)$ depend on $q$ cells varying with $v$, so all $qM$ must be revealed
before the test is fixed; and Lemma~\ref{lem:mass} is then tight for what it is
given, because coordinatewise $\delta$-unpredictability does not yield
$\Exp[m_1m_2]\le\delta^{2}$ (Remark~\ref{rem:corr}). An observer whose query
positions are challenge-independent only \emph{in law}, steering once its coins
are fixed, has per-coin resolution above $1$ and sits here rather than in
Theorem~\ref{thm:cprime}'s first branch. Two exits, neither taken: replace the
reference $p_\theta$ by $\theta(f'(\vx))$ for an independent copy $f'$ of $f$ off
a single cell, which would cut the budget from $qM$ to $q$ at the price of
proving Lemma~\ref{lem:disc} for a random, $f$-dependent reference; or bound
$\Pi(\mathcal S)$ for \emph{reachable} query sets (unions of $M$
decision-tree paths) rather than arbitrary $qM$-sets.

\paragraph{Barrier 2: large $P$ with growing $q$.} Theorem~\ref{thm:C} loses
$P\delta$ at $\mathsf G_2\to\mathsf G_3$, capping the range at
$P\le\sqrt{\sigma'/\delta}$; the statement document names this as a separate open
problem and nothing here improves it. The slack is identifiable: on
$\{\vx\in I_J\}$, which is what costs $P\delta$, consistency makes the decomposed
challenge $H^{*}(\vx)$ equal the real challenge $H(\vx)$, so the event does not
harm $\Adv_{\mathcal Y,D}$ at all, only the route through the uniform-challenge
worlds. A proof comparing $\mathsf G_3$ with $\mathsf G_0$ directly would not pay
it.

\section{External results used}

\begin{itemize}
\item \textbf{[CDGS, Claim~2]} (Coretti, Dodis, Guo, Steinberger, \emph{Random
  Oracles and Non-Uniformity}, ePrint 2017/937, App.~A). \emph{A distribution on
  $[M]^{n}$ with min-entropy deficiency $S$ is $\gamma$-close to a convex
  combination of finitely many $(P',1-\delta)$-dense sources,
  $P'=(S+\log1/\gamma)/(\delta\log M)$.} The paper states this for a uniform
  variable conditioned on a deterministic function; its proof uses only the
  min-entropy deficiency, and the form above is what that proof establishes. The
  proof's structure (conditioning on $\{Y_I=y_I\}$ and recursing on
  $\{Y_I\ne y_I\}$) also supplies the disjointness of the components'
  supports, the fact that fixed values are values the sampled $f$ takes, and,
  via its recursion invariant, that the decomposition is of the conditioned
  distribution itself, so each component is its restriction and the residue's
  mass is at most $\gamma$ measured under it and not under the uniform
  distribution. All three are used in Lemma~\ref{lem:P}.
\item \textbf{[CDGS, Claim~3]} (same paper). \emph{For a $(P',1-\delta)$-dense
  source $X'$ and its corresponding $P'$-bit-fixing $Y'$, and any adaptive
  $T$-query distinguisher, $|\Prob{D^{X'}=1}-\Prob{D^{Y'}=1}|\le T\delta\log M$.}
\item \textbf{[CFHS, Lemma~3]} (Coretti, Farshim, Harasser, Southern,
  \emph{Multi-Source Randomness Extraction and Generation in the ROM}, ePrint
  2025/1258). \emph{For $k\in\mathbb R_{>0}$ with $2^{k}\in\mathbb N$, every
  $k$-source is a convex combination of flat $k$-sources.}
\item \textbf{[CFHS, Lemma~4.3]} \emph{$(n/k)^{k}\le\binom nk\le(en/k)^{k}$.}
\item \textbf{[CFHS, Theorem~3]} \emph{For $\ell$ sources each
  $(N^{\ell},2^{-k})$-unpredictable with unbounded oracle access and a
  $q_D$-query distinguisher,
  $\Adv^{\mathrm{mse}}=O(\ell\log M\sqrt[\ell+1]{q_D2^{-k\ell}(\sigma+\ell N)})$.}
  Used only for the compression arm.
\item \textbf{Hoeffding's inequality}, two-sided, for independent variables with
  ranges of length $c_i$.
\end{itemize}

\section{Review status}\label{sec:review}

This document has \textbf{not} been independently verified, and a reader should
weigh it accordingly.

Two rounds of adversarial review were run, six passes recording a verdict and two
aborting on infrastructure. \textbf{No pass returned clean.} Five findings were
upheld across the two rounds, of which two were serious; both were repaired, and
both were of the same species: a claim left underdetermined rather than false,
each recoverable from machinery already present. Zero numbered conclusions were
falsified at any point.

Three limitations of that review are material. Every referee, the adjudicator and
the author are the same vendor's models, so the errors that survive are the
correlated ones by construction; both serious findings came from the strongest
model and were missed by weaker ones examining the same lines. No pass satisfied
the campaign's own blindness condition, which requires a process carrying none of
the project context and was unavailable in this environment. And no protocol
document existed, so every referee ran on instructions written by the session
that produced the proof.

The text above is the third revision. Its own tally is empty: nothing in it has
been reviewed, and the repair that produced it touched
Definition~\ref{def:res}, on which every resolution argument depends.

\end{document}
